\emph{Exact spectrum structure of norm-equalized coincident copies at general copy number.} Denote by \(T\) the \(2q\times2q\) pencil of the paragraph above, set \(H=\operatorname{Im}C_0=(C_0-C_0^\dagger)/(2\mathrm i)\) and \(K=Z_0-C_0^\dagger C_0=R^\dagger(\mathbf 1_{qN}-UU^\dagger)R\), positive semidefinite. For every \(q\) the characteristic polynomial of the pencil is exactly \(\det(\lambda\mathbf 1_{2q}-T)=\det(\lambda^{2}\mathbf 1_q+4\lambda H-4K)\), so the nonzero spectrum of \(\widetilde Q\) depends on the edge block only through \((H,K)\) and \(\operatorname{Re}C_0\) is invisible at every copy number, not only at \(q=2\) (exact; symbolic at \(q=2\) and \(q=3\), and at \(q=2\) with \(\operatorname{Re}C_0\) entering \(C_0\) and \(Z_0\) jointly; numerical to \(10^{-14}\) through \(q=6\)). Under exact norm equalization, \(K=n\,\mathbf 1_q\) with n>0, which in the block-diagonal per-copy reading is equality of all \(\|w_\perp\|^{2}\) with vanishing cross terms, the \(2q\) nonzero eigenvalues are exactly \(-2h_f\pm2(h_f^{2}+n)^{1/2}\), with \(h_{1},\dots,h_{q}\) the eigenvalues of \(H\), for every \(q\) and with no commutativity hypothesis on \(H\): the Row length pair formula, lifted from \(q=2\) to general \(q\) conditionally on exact equalization, the reduction the paragraph above left unestablished for \(q>2\) (exact; the \(\widetilde Q\)-spectra agree with the pairs to \(2\times10^{-14}\) through \(q=6\)). Within the equalized class a double eigenvalue occurs if and only if two imaginary parts coincide, \(h_f=h_g\), and the multiset is then tiled by the repeated Row length pair of the common value; imaginary-part mismatch leaves the \(2q\) eigenvalues pairwise distinct, with spacing at least \(\min\bigl(2(1-M(n+M^{2})^{-1/2})\min_{f\ne g}|h_f-h_g|,\;4n((n+M^{2})^{1/2}+M)^{-1}\bigr)\) for \(|h_f|\le M\), the gap turning on linearly in the mismatch with log-log slope between \(0.987\) and \(1\) over three decades, and the double-pair locus is the eigenvalue-coincidence set of the Hermitian matrix \(H\), of real codimension three in the \(q^{2}+1\) real parameters \((H,n)\): the only structure forced at mismatch is the fixed-point-free involution \(\lambda\mapsto-4n/\lambda\), which swaps the two members of each Row length pair without ever identifying them, and multiplicity two is forced nowhere except at the excluded boundary \(K=0\), where the pencil loses rank. Beyond exact equalization the quartic at \(q=2\) has a double root exactly where its discriminant vanishes - one real polynomial condition in the eight real parameters of \((H,K)\), and in the diagonal directions tested below it cuts a hypersurface of codimension exactly one - and this double-root locus is strictly larger than the double-pair locus of the paragraph above: the mismatched rows \((\operatorname{Im}\alpha,\|w_\perp\|^{2})=(3,7)\) and \((5,11)\) share the root \(\lambda=2\) with quartic exactly \((\lambda-2)^{2}(\lambda+14)(\lambda+22)\), a double root that is not a double pair, splitting linearly, minimum gap \(1.67\times10^{-4}\) at the shift \(10^{-3}\) of the second norm (exact in rational arithmetic for the factorization, numerical for the split; \(400/400\) random general \((H,K)\) samples have simple spectrum). These are exact-spectrum structure statements for the rank-one edge-datum family of coincident copies at fixed \(q\) and arbitrary \(N\), inside the transplanted rule with its single-dominant-diagram hypothesis at leading semiclassical order and its diagonal-cut convention, valid in the exactly equalized class and, for the last display, in the general \(q=2\) class; they are statements about pencil spectra, not about any state or ensemble, they leave the two-sided reading of the \(\pm\) pairs, the selection rule itself and the general-\(q\) criterion of the paragraph above untouched outside the equalized class, and they close nothing in Problem~G.
